This study evaluated Redis Streams 7.4.2 and NATS JetStream 2.10.24 through explicit protocol analysis, conditional mathematical proofs, selected machine-checked model results, and retained homelab measurements. Its conclusions are bounded by single-replica deployments, the specified client paths, compressible 512-byte payloads, and one shared storage environment.

For RQ1, observed failure-to-redelivery intervals were consistent with the remaining eligibility timeout and subsequent recovery scheduling. Later kills were followed by shorter intervals in both systems, and longer Redis reclaim sleeps were associated with longer additional delays. All ten new-read-only controls remained pending. All 120 active recovery episodes completed with the original identifier and cleared pending state. The residual-time theorem explains this pattern under stated timing assumptions. Neither the observations nor the theorem establish death detection or guaranteed empirical deadlines.

For RQ2, mean acknowledged backlog-drain throughput with sixteen logical consumers was 8.59--12.94 times the one-consumer mean across profiles. Cycle p99 was higher while the backlog completed sooner. The verified finite-window occupancy conditions and the identity $X=CU/\overline R$ connect this sublinear scaling to the measured cycle cost. They do not establish an unbounded scaling law, and the different always-mode consumer-state persistence paths prevent interpreting that cross-system drain comparison as equal-guarantee performance.

For RQ3, periodic configurations had approximately 83\% of Redis memory-mode publication throughput and 82\% of JetStream memory-mode throughput. Always-mode publication exhibited lower throughput and greater variability, including one failed planned Redis trial. Completed-run means and nominal intervals retain that conditioning; the failure and operational deviations remain visible. Leave-one-block-out diagnostics preserve the directions of paired contrasts, while influential observations, sparse tail information, execution-order variation, and unequal storage-history carryover limit precision and prevent isolating synchronization-policy effects. The separate serial sensitivity further shows why duration and request concurrency must accompany a throughput figure. No consumer-kill or publication-latency result establishes power-loss durability.

The contribution is an auditable, computationally reproducible comparison of specific recovery policies, concurrency behavior, and acknowledgment costs under documented conditions. The versioned, licensed artifact traces numerical claims to retained observations and formal claims to their premises. Future work should repeat the study across independent storage hosts and less-compressible payloads, hold duration fixed while varying publication concurrency, and add replication, large pending sets, open-loop load, and controlled broker/storage faults. These extensions would address external validity and failure classes that the present evidence leaves unresolved.
